\documentclass[11pt]{article}

\usepackage[margin=1in]{geometry}
\usepackage{amsmath,amssymb,amsthm,mathtools}
\usepackage{booktabs}
\usepackage{enumitem}
\usepackage{hyperref}
\usepackage{xcolor}
\hypersetup{hidelinks}

\newtheorem{theorem}{Theorem}[section]
\newtheorem{lemma}[theorem]{Lemma}
\newtheorem{proposition}[theorem]{Proposition}
\newtheorem{corollary}[theorem]{Corollary}
\newtheorem{observation}[theorem]{Observation}
\theoremstyle{definition}
\newtheorem{definition}[theorem]{Definition}
\theoremstyle{remark}
\newtheorem{remark}[theorem]{Remark}

\newcommand{\calG}{\mathcal{G}}
\newcommand{\calP}{\mathcal{P}}
\newcommand{\calC}{\mathcal{C}}
\newcommand{\calB}{\mathcal{B}}
\newcommand{\calS}{\mathcal{S}}
\newcommand{\calK}{\mathcal{K}}
\newcommand{\LL}{\mathrm{LL}}
\newcommand{\FC}{\mathrm{FC}}
\newcommand{\rev}{\operatorname{rev}}

\title{Zero forcing on temporal graphs:\\
Layer-local and footprint-constrained dynamics}
\author{Krishna Harish\\
Elkins High School, Missouri City, Texas, USA}
\date{}

\begin{document}
\maketitle

\begin{abstract}
Zero forcing is a graph-coloring process in which a blue vertex with a unique
white neighbor forces that neighbor to blue.  The standard process assumes a
fixed graph, whereas many communication, interaction, and control networks
change over time.  We study two synchronous finite-lifetime extensions to a
temporal graph.  In \emph{layer-local temporal zero forcing}, uniqueness is
evaluated in the current layer.  In \emph{footprint-constrained temporal zero
forcing}, uniqueness is evaluated in the aggregate footprint while the
forcing edge must be currently active.  The distinction changes the process
fundamentally: layer-local forces may branch over time, while
footprint-constrained forces form vertex-disjoint chains.

We prove sharp universal growth bounds, dominance and reversal results, and
exact formulas for alternating temporal paths and even cycles.  These formulas
give linear separations between the two temporal parameters and ordinary zero
forcing of the aggregate graph.  A temporal-tree family shows that static
aggregation can overestimate layer-local forcing by a linear factor and that
reversing the layers can change the layer-local number by a factor linear in
the order; the footprint-constrained number is invariant under reversal.  We
also prove that layer-local temporal zero forcing is NP-complete and
W[2]-hard by seed budget even when every layer is a star, every edge appears
once, and the footprint is bipartite.  Finally, the growth bounds yield
fixed-parameter algorithms parameterized by seed budget and lifetime.
Exhaustive computations verify all closed forms on small instances.
\end{abstract}

\section{Introduction}

Zero forcing was introduced as a combinatorial tool for minimum-rank
problems and has subsequently been connected with graph searching,
propagation, and network controllability~\cite{aim2008,yang2013,fallat2016,
monshizadeh2014}.  Starting
from an initial set of blue vertices, the color-change rule repeatedly allows
a blue vertex with exactly one white neighbor to color that neighbor blue.
The minimum size of an initial set that eventually colors the graph is the
zero-forcing number.

The static graph assumption is restrictive when edges represent contacts or
communication opportunities.  Temporal-graph models retain the order in
which edges are available, and temporal order can change reachability,
influence, and controllability even when the aggregate graph is unchanged
~\cite{casteigts2012}.
Existing work has temporalized target-set selection~\cite{schierreich2023},
compared temporal and aggregate influence~\cite{michalski2014}, minimized
temporal source sets~\cite{yadav2026}, and studied controllability of
time-varying networks~\cite{srighakollapu2022,mousavi2021,yao2017}.
Applying the unique-white-neighbor rule to temporal layers introduces an
additional modeling choice: should
uniqueness be evaluated among currently active neighbors or among all
neighbors that can ever appear?

We formalize both choices.  The layer-local rule evaluates uniqueness in the
current layer.  It models a source whose ambiguity is limited to its current
contacts.  The footprint-constrained rule evaluates uniqueness in the
aggregate footprint and then requires the unique remaining edge to be active.
It models a temporally scheduled version of a valid static zero-forcing
deduction.  The two rules agree on neither their growth mechanism nor their
response to aggregation and time reversal.

Our contributions are as follows.

\begin{enumerate}[leftmargin=*]
  \item We prove
  \[
    Z_{\LL}(\calG)\le Z_{\FC}(\calG),\qquad
    Z_{\LL}(\calG)\ge \left\lceil\frac{n}{2^T}\right\rceil,
    \qquad
    Z_{\FC}(\calG)\ge
    \max\left\{Z(H),\left\lceil\frac{n}{T+1}\right\rceil\right\},
  \]
  where $T$ is the lifetime and $H$ is the footprint.  Both growth bounds
  are sharp.
  \item For the two-layer alternating path $\calP_n$ and even cycle
  $\calC_n$, we determine both parameters exactly:
  \[
  Z_{\LL}(\calP_n)=\left\lfloor\frac n4\right\rfloor+1,
  \quad Z_{\FC}(\calP_n)=\left\lceil\frac n3\right\rceil,
  \]
  and
  \[
  Z_{\LL}(\calC_n)=\left\lceil\frac n4\right\rceil,
  \quad Z_{\FC}(\calC_n)=2\left\lceil\frac n6\right\rceil.
  \]
  \item We construct a temporal tree on $n=2^T$ vertices for which the
  forward layer-local number is $1$, the footprint-constrained number is
  $n/2$, and the layer-local number after reversing time is $n/2$.
  Meanwhile, the static zero-forcing number of the footprint is $n/4$.
  Thus the family gives linear additive aggregation errors in both
  directions.  Together with the alternating paths, it shows that either
  multiplicative error can grow linearly with $n$.
  \item We prove footprint-constrained forcing is invariant under reversing
  the layer order, in contrast to the unbounded directional asymmetry of the
  layer-local rule.
  \item We establish restricted NP-completeness and W[2]-hardness for the
  layer-local problem through an additive optimum-preserving reduction from Set
  Cover, and derive fixed-parameter algorithms from the two growth laws.
\end{enumerate}

\section{Related work and prior-art boundary}

Standard zero forcing and its minimum-rank origin are developed in
\cite{aim2008}.  Parallel rounds and propagation time on a fixed graph are
studied in~\cite{hogben2012}, and recent work has introduced relaxed
chronologies and parallel increasing path covers~\cite{hogben2025}.  We
therefore avoid using ``chronological zero forcing'' as the name of the new
parameters: chronology is already established terminology in the static
literature.

A closely related temporal diffusion model is temporal target-set selection.  Its
update is centered on the target: a vertex activates when sufficiently many
of its current neighbors are active~\cite{schierreich2023}.  Our condition is
centered on the source: a blue source can act only when it has a unique white
neighbor.  Minimum temporal source sets based on time-respecting
reachability~\cite{yadav2026} likewise do not depend on the colors of a
source's other neighbors.

There is also a close structural boundary with temporal path covers.  Every
footprint-constrained process, after one parent is selected for each newly
blue vertex, decomposes into vertex-disjoint strict temporal paths.  Temporal
path-cover problems optimize path collections subject to temporal
reachability or disjointness~\cite{chakraborty2024,cioni2026}; an arbitrary
temporal path cover need not satisfy the footprint-uniqueness condition that
makes a path a valid forcing chain.  The footprint-constrained parameter can
thus be viewed as a uniqueness-constrained temporal path-cover number, not a reformulation of the
unconstrained path-cover problems.

Temporal controllability is a mature adjacent literature.  In particular,
zero forcing has already been used for time-varying structural
perturbations~\cite{mousavi2021}, and direct graph conditions and input-set
optimization are known for temporal strong structural
controllability~\cite{srighakollapu2022}.  We make no claim here that either
temporal color process characterizes a previously studied matrix class.

The one-layer boundary is also known under different language.  If $W$ is
the set left white initially, one synchronous layer succeeds exactly when
every $w\in W$ has an external private neighbor.  Writing
$\operatorname{OIR}(G)$ for the maximum cardinality of an open irredundant
set, this gives the exact identity
\[
 Z_{\LL}((V,E_1))=|V|-\operatorname{OIR}(G_1).
\]
Open irredundance and its upper parameter are studied in
\cite{fellows1994}.  Similarly, integer programming for bounded-time static
zero forcing predates the present temporal model~\cite{brimkov2019}.  Our
complexity contribution therefore uses a multi-layer temporal sequence in
which each edge appears exactly once.

\section{Definitions}

All graphs are finite, simple, and undirected.  A temporal graph is an ordered
tuple
\[
  \calG=(V,E_1,\ldots,E_T),
\]
where $G_t=(V,E_t)$ is the layer at time $t$.  The lifetime is $T$, and the
\emph{footprint} is
\[
  H=\left(V,\bigcup_{t=1}^T E_t\right).
\]
Colors persist.  If $B\subseteq V$ is the blue set immediately before layer
$t$, define
\begin{align*}
F_t^{\LL}(B)
 &=\left\{v\notin B:
 \exists u\in B, uv\in E_t,
 N_{G_t}(u)\setminus B=\{v\}\right\},\\
F_t^{\FC}(B)
 &=\left\{v\notin B:
 \exists u\in B, uv\in E_t,
 N_H(u)\setminus B=\{v\}\right\}.
\end{align*}
For $X\in\{\LL,\FC\}$, the synchronous process from a seed set $S$ is
\[
 B_0^X=S,
 \qquad B_t^X=B_{t-1}^X\cup F_t^X(B_{t-1}^X).
\]
Eligibility is always evaluated using $B_{t-1}^X$; newly blue vertices cannot
force until the next layer.

\begin{definition}
A set $S\subseteq V$ is an $X$-temporal zero-forcing set if $B_T^X=V$.
The minimum cardinality of such a set is $Z_X(\calG)$.
\end{definition}

If several sources can force the same target in one round, the target is
colored once.  When a parent relation is needed, we choose one valid source.

\section{General structure and sharp growth bounds}

\begin{proposition}[Dominance]\label{prop:dominance}
For every temporal graph $\calG$,
\[
 Z_{\LL}(\calG)\le Z_{\FC}(\calG).
\]
\end{proposition}

\begin{proof}
Run both processes from the same seed set.  We prove inductively that
$B_t^{\FC}\subseteq B_t^{\LL}$.  Suppose $u$ footprint-forces $v$ at time
$t$.  Every footprint neighbor of $u$ other than $v$ is already blue in the
footprint process and hence, by induction, in the layer-local process.  If
$v$ is not already layer-locally blue, it is the unique white neighbor of
$u$ in $G_t$ and is forced there as well.  Thus every footprint-constrained
forcing set is layer-local.
\end{proof}

\begin{theorem}[Growth dichotomy]\label{thm:growth}
If $|V|=n$ and $\calG$ has lifetime $T$, then
\[
 Z_{\LL}(\calG)\ge \left\lceil\frac{n}{2^T}\right\rceil
\]
and
\[
 Z_{\FC}(\calG)\ge
 \max\left\{Z(H),\left\lceil\frac{n}{T+1}\right\rceil\right\}.
\]
Both lifetime-dependent bounds are attained for infinitely many orders.
\end{theorem}

\begin{proof}
In a layer-local round each blue vertex has at most one unique white
neighbor, so at most $|B_{t-1}^{\LL}|$ vertices are added.  Hence
$|B_t^{\LL}|\le 2|B_{t-1}^{\LL}|$ and
$n\le 2^T|S|$.

In the footprint-constrained process, a vertex can force at most once over
the entire lifetime.  Immediately after it forces its unique footprint-white
neighbor, every one of its footprint neighbors is blue permanently.  After
choosing one parent for every newly blue vertex, the forces therefore form
vertex-disjoint directed paths rooted at the seeds.  Time labels increase
strictly along each path, so a path contains at most $T+1$ vertices.  Thus
$n\le (T+1)|S|$.  Moreover, every footprint-constrained force is a valid
ordinary zero-forcing move in $H$, so $S$ is a static zero-forcing set and
$|S|\ge Z(H)$.

The layer-local bound is attained by the binomial temporal trees in
Section~\ref{sec:tree}.  The footprint bound is attained by disjoint
$(T+1)$-vertex paths whose successive edges appear at times $1,\ldots,T$.
\end{proof}

\section{Exact alternating paths and cycles}\label{sec:exact}

Let $P_n$ have vertices $1,\ldots,n$ in order.  Define the two-layer temporal
path $\calP_n$ by
\[
 E_1=\{(1,2),(3,4),\ldots\},
 \qquad
 E_2=\{(2,3),(4,5),\ldots\}.
\]

\begin{theorem}[Alternating paths]\label{thm:path}
For every $n\ge1$,
\[
 Z_{\LL}(\calP_n)=\left\lfloor\frac n4\right\rfloor+1,
 \qquad
 Z_{\FC}(\calP_n)=\left\lceil\frac n3\right\rceil.
\]
\end{theorem}

\begin{proof}
For the layer-local rule, let
$Q_k=\{2k-1,2k\}\cap[n]$.  A seed in $Q_k$ makes its first-layer pair blue
and can then expand across both adjacent second-layer edges.  Its complete
reach is
\[
 I_k=\{2k-2,2k-1,2k,2k+1\}\cap[n].
\]
Because both layers are matchings, multiple seeds color exactly the union of
their intervals $I_k$.  Each interval contains at most four vertices, so
$\lceil n/4\rceil$ seeds are necessary.  When $n=4a$, vertex $1$ belongs
only to the three-vertex interval $I_1$; consequently $a$ intervals cover at
most $3+4(a-1)=n-1$ vertices.  This handles $n\equiv0\pmod4$; combined with
$\lceil n/4\rceil$ for the other residues, it gives
$\lfloor n/4\rfloor+1$ in every case.  For the upper bound, write
$n=4a+r$.  If $r=0$, use the interval indices
$\{1,3,\ldots,2a-1,2a\}$; if $r\in\{1,2,3\}$, use
$\{1,3,\ldots,2a+1\}$.  (For $a=0$, the latter set is $\{1\}$.)  These
intervals cover $[n]$ with exactly
$\lfloor n/4\rfloor+1$ choices; seed any vertex of the corresponding $Q_k$.

For the footprint-constrained rule, Theorem~\ref{thm:growth} with $T=2$
gives the lower bound $\lceil n/3\rceil$; the cases $n\le2$ are direct.  For
the upper bound, partition the first $3q$ vertices into triples
$D_j=\{3j-2,3j-1,3j\}$.  If $j$ is odd, seed $3j-2$; if $j$ is even, seed
$3j$.  For an internal odd block, the left neighbor outside the block is the
root of the preceding even block and is already blue; hence the odd root has
its middle vertex as unique white footprint neighbor and forces rightward in
layers $1$ and $2$.  Symmetrically, an internal even root is supported by the
root of the following odd block and forces leftward.  Endpoint blocks need no
outside support.  Thus $q$ seeds color the first $3q$ vertices when
$n=3q$.  If $n=3q+1$, seed $n$; this also supports a final even block.  If
$n=3q+2$ and $q$ is even, seed $n-1$: it supports the final even root and
forces $n$ in layer $1$.  If $q$ is odd, seed $n$, which forces $n-1$ in
layer $2$.  These constructions use exactly $\lceil n/3\rceil$ seeds.
\end{proof}

Since the footprint of $\calP_n$ is $P_n$ and $Z(P_n)=1$, both temporal
parameters exceed their aggregate static counterpart by a linear factor.
The two temporal semantics themselves differ by $n/12+O(1)$.

For even $n\ge4$, define $\calC_n$ by placing alternating perfect matchings
of $C_n$ in its two layers.

\begin{theorem}[Alternating even cycles]\label{thm:cycle}
For even $n\ge4$,
\[
 Z_{\LL}(\calC_n)=\left\lceil\frac n4\right\rceil,
 \qquad
 Z_{\FC}(\calC_n)=2\left\lceil\frac n6\right\rceil.
\]
\end{theorem}

\begin{proof}
The case $n=4$ follows directly, so assume $n\ge6$.  For layer-local forcing,
each seed in a first-layer pair reaches a cyclic arc
of four vertices.  Hence at least $\lceil n/4\rceil$ seeds are necessary.
Selecting alternating first-layer pairs covers the cycle with exactly that
many arcs.

For the footprint-constrained upper bound, regard the $n/2$ second-layer
edges as cyclic blocks.  Seed both endpoints of every block in a minimum
dominating set of this block cycle.  In round one, the endpoints of a selected
block force their outside first-layer neighbors.  In round two, those newly
blue vertices force the as-yet-uncolored endpoints of the two adjacent
second-layer blocks.  Thus a selected block covers itself and its two
neighboring blocks, using
$2\lceil(n/2)/3\rceil=2\lceil n/6\rceil$ seeds.

The chain bound gives $n\le3|S|$.  It remains only to exclude the odd value
$|S|=2q+1$ when $n=6q+2$.  Let $f$ be the number of distinct vertices added
in round one.  Since $E_1$ is a matching, $f$ is also the number of seeds
that force in that round.  If
$f\le |S|-2$, then before round two there are $|S|+f$ blue vertices, but the
$f$ first-round sources can never force again under the footprint-constrained
rule.  Hence at most $|S|$ vertices can force in round two, and
$n\le |S|+f+|S|\le3|S|-2$.  Equality $f=|S|$ would require every seed to
have exactly one blue cycle neighbor, so the seed-induced subgraph would be
a disjoint union of edges and $|S|$ would be even.  If $f=|S|-1$, let $x$ be
the unique first-round nonforcer.  If $x$ has no blue cycle neighbor, its
$E_1$ neighbor cannot become blue in round one and its only possible
unique-white edge is unavailable in round two.  If $x$ has two blue
neighbors, it has no white neighbor.  If it has one blue neighbor, then every
seed has degree one in the seed-induced subgraph, which would be 1-regular
and would force $|S|$ to be even.  Hence $x$ cannot force in round two, so the
second round adds at most $|S|-1$ vertices and again
$n\le3|S|-2$.  Thus $n=6q+2$ requires at least $2q+2$ seeds, completing all
residue classes.
\end{proof}

\section{Aggregation, semantics, and time reversal}\label{sec:tree}

We next give one family witnessing three distinct failures of a static-only
view.  Start with $V_0=\{r\}$.  At time $t$, create one fresh child
$u^{(t)}$ for every $u\in V_{t-1}$, put
\[
 L_t=\{u^{(t)}:u\in V_{t-1}\},\quad
 E_t=\{uu^{(t)}:u\in V_{t-1}\},\quad
 V_t=V_{t-1}\cup L_t.
\]
Let $\calB_T$ be the resulting temporal graph and $H_T$ its footprint.  Each
layer is a matching, $H_T$ is a tree, and $|V_T|=2^T=n$.

\begin{theorem}[Binomial temporal tree]\label{thm:binomial}
For $T\ge2$,
\[
 Z_{\LL}(\calB_T)=1,
 \qquad
 Z_{\FC}(\calB_T)=\frac n2,
 \qquad
 Z(H_T)=\frac n4.
\]
If the layers are reversed, then
\[
 Z_{\LL}(\rev(\calB_T))=\frac n2.
\]
\end{theorem}

\begin{proof}
Seed $r$.  Inductively, all of $V_{t-1}$ is blue before layer $t$, and every
blue vertex has exactly one current white neighbor, its new child.  Thus
$Z_{\LL}(\calB_T)=1$.

In the first layer of the reversed process, $E_T$ is a perfect matching
between $V_{T-1}$ and the leaves $L_T$.  If an edge has neither endpoint
seeded, its leaf can never be colored because the edge never reappears.
Therefore every reversed layer-local forcing set has at least
$|E_T|=n/2$ vertices.  Seeding all leaves $L_T$ colors their parents in the
first reversed round and hence colors the whole graph.  The same leaf set is
valid footprint-constrained, while Proposition~\ref{prop:dominance} supplies
the lower bound.  For the footprint-constrained forward lower bound, every
forcing chain contains at most one vertex of $L_T$: the sole incident edge of
such a leaf has label $T$, while labels increase strictly along a chain.
Since $|L_T|=n/2$, at least $n/2$ chains and hence seeds are required.
Directly, seeding $V_{T-1}$ lets every old vertex force its time-$T$ child,
proving equality.

The leaves of $H_T$ are exactly $L_T$, so there are $n/2$ of them.  Static
forcing chains form vertex-disjoint paths, each containing at most two tree
leaves, and therefore $Z(H_T)\ge n/4$.  For the reverse inequality, seed the
time-$T$ child of every vertex in $V_{T-2}$.  These leaves force
$V_{T-2}$; those vertices then force their time-$(T-1)$ children; finally the
time-$(T-1)$ children force their own time-$T$ children.  The seed set has
$|V_{T-2}|=n/4$ vertices.
\end{proof}

Together with Theorem~\ref{thm:path}, the family shows that static
aggregation can either underestimate or overestimate layer-local temporal
zero forcing by $\Theta(n)$.

An even simpler family gives the largest possible order of overestimation up
to an additive constant.

\begin{proposition}[Serialized stars]\label{prop:star}
Let $\calS_n$ be the temporal star with center $c$, leaves
$v_1,\ldots,v_{n-1}$, and one-edge layers $E_t=\{cv_t\}$ in this order.
For $n\ge3$,
\[
 Z_{\LL}(\calS_n)=1,
 \qquad
 Z_{\FC}(\calS_n)=Z(K_{1,n-1})=n-2.
\]
\end{proposition}

\begin{proof}
For the layer-local rule, seed the center; it forces one leaf in every layer.
The footprint-constrained number is at least the static number
$Z(K_{1,n-1})=n-2$: at most one leaf can be omitted from a static forcing
set, because two white leaves leave the center unable to distinguish either;
seeding all but one leaf is sufficient.  For temporal equality, seed
$v_1,\ldots,v_{n-2}$.  Vertex $v_1$ forces $c$ in the first layer, and at the
last layer $c$ has only $v_{n-1}$ white and forces it.
\end{proof}

\begin{corollary}[Maximum additive aggregation overestimate]\label{cor:maxgap}
For every $n\ge3$ there is a temporal graph $\calK_n$ with
\[
 Z_{\LL}(\calK_n)=1,
 \qquad
 Z_{\FC}(\calK_n)=Z(K_n)=n-1.
\]
Consequently, the additive overestimate
$Z(H)-Z_{\LL}(\calG)$ can attain its maximum possible value $n-2$.
\end{corollary}

\begin{proof}
First serialize the $n-1$ edges from a center $c$ to the other vertices, and
then place every remaining edge of $K_n$ in its own later layer.  Seeding $c$
colors every other vertex during the initial star layers, so the later layers
are irrelevant to layer-local forcing.  The footprint is $K_n$, whose static
zero-forcing number is $n-1$.  The footprint lower bound gives
$Z_{\FC}\ge n-1$.  For the reverse inequality, seed every vertex except the
last star leaf.  At its star layer, $c$ has that leaf as its unique white
footprint neighbor and the corresponding edge is active, so $c$ forces it.
Finally, if a footprint has an edge then $Z(H)\le n-1$ and
$Z_{\LL}\ge1$; for an edgeless footprint both numbers equal $n$.
\end{proof}

Static zero-forcing chains can be reversed to obtain a zero-forcing set of
the same size~\cite{barioli2010}.  The next result shows that this operation
survives the temporal schedule under footprint-constrained forcing.

\begin{theorem}[Footprint-constrained reversal]\label{thm:reversal}
For every temporal graph,
\[
 Z_{\FC}(\calG)=Z_{\FC}(\rev(\calG)).
\]
\end{theorem}

\begin{proof}
Take a successful forward process and choose one parent force for every newly
blue vertex, discarding duplicate forces to the same target.  By footprint
uniqueness, every vertex has at most one selected outgoing force; every
nonseed has one selected incoming force.  The selected arcs are therefore
vertex-disjoint directed paths or directed cycles.  Their time labels increase strictly along
each directed walk because a newly colored vertex cannot force in the same
round.  In particular, there is no directed cycle, so all components are
paths.  Every path has exactly
one original seed and one terminal vertex, so the number of terminals equals
the number of seeds.

Seed the reversed process with the terminal vertices of these paths and
reverse every arc $x\to y$ used at time $t$ to $y\to x$ at time
$T+1-t$.  We verify this schedule by induction over decreasing forward time
labels.  Vertex $y$ is already blue then, either because it is terminal or
because its later forward outgoing arc has already been reversed.  Let
$w\in N_H(y)\setminus\{x\}$.  If $w$ has no selected outgoing force, it is a
terminal vertex and hence a seed of the reversed process.  Otherwise let
$w\to z$ occur at forward time $s$.  If $s<t$, then $y$ was white before
time $s$; footprint uniqueness forces $z=y$, so $y$ would become blue at
time $s$, contradicting the selected force $x\to y$ at time $t$.  If $s=t$,
the source $w$ can have only the unique target $y$; after selecting $x\to y$
as the parent of $y$, the duplicate $w\to y$ is not selected, contradicting
the assumption that $w$ has a selected outgoing force.  Hence $s>t$, so
$z\to w$ has already
occurred in the prescribed reverse schedule.  Thus every footprint neighbor
of $y$ other than $x$ is blue.  At reverse time $T+1-t$, either $x$ has
already become blue through an additional deterministic force, or $x$ is the
unique white footprint neighbor of $y$ and the active edge $xy$ forces it.
By blue-set monotonicity, additional forces can only help; all prescribed
targets are therefore blue by their scheduled reverse times.  This proves
$Z_{\FC}(\rev(\calG))\le Z_{\FC}(\calG)$; reversing twice gives equality.
\end{proof}

Theorem~\ref{thm:binomial} shows that no analogous result holds for the
layer-local rule: its reversal ratio can be $n/2$.

\section{Computational complexity and parameterized algorithms}

The decision problem \textnormal{\textsc{LL-TZF}} asks whether
$Z_{\LL}(\calG)\le K$.  Membership in NP follows by simulating the
deterministic process.

\begin{theorem}[Restricted hardness]\label{thm:hardness}
\textnormal{\textsc{LL-TZF}} is NP-complete even when every layer has one nontrivial
component that is a star, every footprint edge appears in exactly one layer,
and the footprint is bipartite.  Parameterized by $K$, it is W[2]-hard.
The reduction satisfies
\[
 \operatorname{OPT}_{\textnormal{\textsc{LL-TZF}}}
 =\operatorname{OPT}_{\textnormal{\textsc{Set Cover}}}+1.
\]
\end{theorem}

\begin{proof}
Let a Set Cover instance have universe
$U=\{e_1,\ldots,e_p\}$ and sets $\mathcal F=\{F_1,\ldots,F_m\}$, with
every element occurring in a set.  Create vertices
\[
 \{c,a\}\cup\{s_1,\ldots,s_m\}\cup\{x_1,\ldots,x_p\}.
\]
The first layer is the edge $ca$.  For each element $e_j$, the next layer is
the star centered at $x_j$ with leaves $s_i$ for which $e_j\in F_i$.
Finally, for $i=1,\ldots,m$, add a one-edge cleanup layer $cs_i$.  Every edge
appears once, and the footprint is bipartite with parts
$\{c,x_1,\ldots,x_p\}$ and $\{a,s_1,\ldots,s_m\}$.

If $C$ is a cover, seed $c$ and the set vertices corresponding to $C$.
Vertex $c$ first forces $a$.  At each element layer, a seeded incident set
vertex forces the center $x_j$.  The cleanup layers let $c$ force all
remaining set vertices.  This gives a temporal forcing set of size
$|C|+1$.

Conversely, let $S$ be feasible.  At least one of $c,a$ is seeded, since $a$
has no active edge after the first layer.  Let $P$ index the set vertices
that are blue before cleanup.  Such a set vertex is either seeded or was
forced in an element layer by its center.  A center performing that force
must itself be seeded, because it has no earlier incident edge.  A center
forced during its own element layer cannot force back in that layer under the
synchronous update convention; moreover, any seeded center can force at most
one set vertex.  Therefore the vertices indexed by $P$ can
be injectively charged to seeds outside $\{c,a\}$.

If element $e_j$ is not contained in any set indexed by $P$, then no incident
set vertex was blue at its layer.  Hence $x_j$ must be seeded.  It was not
used in the preceding charge, since forcing an incident set vertex would put
a set containing $e_j$ into $P$.  For every such uncovered element, choose
one arbitrary containing set.  Together with the sets indexed by $P$, these
choices form a cover of size at most $|S|-1$.  Thus the optima differ by
exactly one.  The transformation sends $k$ to $k+1$; parameterized Set Cover
is W[2]-complete~\cite{downeyfellows2013}, establishing W[2]-hardness.
\end{proof}

Using the classical NP-complete Vertex Cover problem~\cite{gareyjohnson1979}
as the incidence instance makes every element star a $P_3$, yielding
NP-completeness even when every layer has at most two edges.

\begin{theorem}[Seed-budget--lifetime algorithms]\label{thm:fpt}
Both temporal zero-forcing problems are fixed-parameter tractable in
$k+T$.  More precisely, after a growth-bound rejection, exhaustive seed
enumeration runs in
\[
 2^{O(k(T+1))}\operatorname{poly}(|I|)
\]
for layer-local forcing and
\[
 2^{O(k\log(T+1))}\operatorname{poly}(|I|)
\]
for footprint-constrained forcing.  The resulting vertex kernels have
$k2^T$ and $k(T+1)$ vertices, respectively; under explicit layer encoding,
the second is a polynomial kernel in $k+T$.
\end{theorem}

\begin{proof}
For a nonempty instance, $k=0$ is immediately a no instance.  Now suppose
$k\ge1$.  First observe that the final blue set is monotone in the seed set.
Indeed,
if $B\subseteq B'$, any force available from $B$ either has its target already
in $B'$ or remains available because enlarging the blue set can only remove
white neighbors.  Induction over the layers proves monotonicity for both
semantics.  Thus a solution of size at most $k<n$ can be extended to one of
size exactly $k$; if $k\ge n$, the instance is trivially feasible.

By Theorem~\ref{thm:growth}, reject a layer-local instance if $n>k2^T$.
Otherwise enumerate the $k$-subsets.  The standard binomial bound
$\binom{n}{k}\le(en/k)^k$ gives $2^{O(k(T+1))}$ candidates.  For the
footprint-constrained problem reject if $n>k(T+1)$; the same estimate gives
$(e(T+1))^k=2^{O(k\log(T+1))}$ candidates.  Each candidate is checked by
direct simulation in polynomial time.  If the size test passes, retaining
the instance gives the stated vertex bounds; otherwise output a fixed no
instance.  With an explicit temporal edge list, $T[k(T+1)]^2$ bounds the
encoding size of the footprint-constrained kernel by a polynomial in $k+T$.
\end{proof}

\section{Computational methods and verification}

We implemented the two synchronous updates directly from their definitions,
without encoding the closed forms, and enumerated seed sets in increasing
cardinality.  The implementation also
computes ordinary static zero forcing by repeated closure.  The checked
instances comprise

\begin{itemize}[leftmargin=*]
  \item alternating paths through $n=20$;
  \item alternating even cycles through $n=22$;
  \item forward and reversed binomial temporal trees through $T=4$
  ($n=16$);
  \item serialized temporal stars through $n=12$; and
  \item completed serialized cliques through $n=8$.
\end{itemize}

All 127 recorded parameter values matched their corresponding predicted
values.  The code
records an optimum witness for every instance and writes machine-readable
CSV and JSON results.  It additionally exhausts all 4,096 two-layer temporal graphs on
four labeled vertices to check dominance and footprint-constrained reversal
invariance (S1 File).  These computations are verification rather than evidence used
in place of proof.

\section{Discussion}

The semantic choice determines whether time-varying neighborhoods create
branching or merely schedule static deductions.  Layer-local forcing can
exploit temporarily absent edges, so it may be dramatically easier than
forcing on the aggregate graph.  The same deadline can also make it
dramatically harder, as the alternating paths demonstrate.
Footprint-constrained forcing is more conservative: it always begins with a
static zero-forcing set, has no branching, and is time-reversal invariant.

These distinctions caution against assigning a single ``temporal zero
forcing number'' without specifying (i) the neighborhood in which uniqueness
is evaluated, (ii) whether layers permit one round or closure, and (iii)
whether simultaneous eligibility uses the pre-round coloring.  Our results
use one synchronous round per layer.  Because the exact path and cycle layers
are matchings, allowing closure within those individual layers would not
change their values, but it changes the general model.

A next structural problem is an explicit dynamic program on temporally
labeled paths and trees beyond the alternating families.  A second direction
is to identify a precisely stated time-varying matrix class for which one of
the parameters provides a sound controllability certificate.  Existing
temporal controllability results mean that such a connection must be proved,
not inferred from the static interpretation.

\section{Conclusion}

Zero forcing on a temporal graph has at least two natural meanings.  Evaluating
uniqueness in the current layer allows repeated branching and strong temporal
asymmetry; evaluating it in the footprint preserves static forcing chains and
time-reversal symmetry.  Exact families, aggregation gaps, complexity, and
parameterized algorithms show that this distinction is mathematical rather
than terminological.  The results provide a starting point for further work on
temporal graph classes and carefully delimited control applications.

\section*{Data and code availability}

All code and machine-readable results underlying the computational verification
are publicly available from Zenodo under the MIT License
(\href{https://doi.org/10.5281/zenodo.21347118}{doi:10.5281/zenodo.21347118})
and from the
\href{https://github.com/krishnatheaverage/temporal-zero-forcing}{development
repository on GitHub}.  The same verification materials are provided as S1
File.  No external datasets were used.

\bibliographystyle{plain}
\bibliography{references}

\section*{Supporting information}

\noindent\textbf{S1 File. Code and machine-readable verification results.}
ZIP archive containing the Python verification script, generated CSV and JSON results,
reproduction instructions, tested Python version, and MIT license.

\end{document}
